\section{Discussion}
\label{sec:discussion}

\subsection{Why SE Wins on TriviaQA}

The +0.155 AUROC gap is explained directly by the mechanism measured in RQ2.
Token entropy aggregates surface variation within semantically equivalent paraphrase groups --- treating different word choices for the same wrong answer as independent uncertainty signals.
Semantic entropy's NLI clustering removes this variation before computing entropy.
The scale of the effect (intra-cluster TE variance = 7.152 nats$^2$, 71$\times$ threshold) makes the explanation unambiguous: paraphrase noise is the dominant component of TE's within-cluster signal.
This extends \citet{Huang2023}'s empirical observation that sampling-based methods outperform single-pass entropy, by providing the feature-level mechanism.

\subsection{Why SE Loses on TruthfulQA}

TruthfulQA's adversarial misconceptions produce low-diversity model outputs: across K=10 samples, Llama-2-7B consistently generates the same wrong answer with high confidence.
All K samples fall into one NLI cluster, yielding near-zero SE for incorrect answers --- making them indistinguishable from correct answers in SE's feature space.
Token entropy succeeds here: a narrow, peaked token distribution (low entropy) for a deterministically wrong output is a reliable incorrectness indicator.

Three competing explanations exist for the reversal: (1) task-structure dependence (primary --- mechanistically consistent with H-M1); (2) NLI model size (H-C1 used \texttt{nli-deberta-v3-small} vs.\ H-E1's large variant --- a genuine confound); (3) EM label noise (low plausibility --- would deflate all AUROCs equally).
We cannot cleanly attribute the reversal to (1) vs.\ (2) without a follow-up experiment holding the NLI model fixed.

\subsection{SCG BERTScore as a Short-QA Failure Mode}

The $|\text{SCG} - \text{SE}|$ gap of 0.336 (corrected) establishes that BERTScore-based consistency is not an SE substitute on short factual QA.
BERTScore measures contextual embedding similarity, not semantic entailment --- on 1--3 word TriviaQA answers, this distinction is critical.
Practitioners seeking lower-compute SE alternatives should use SelfCheckNLI rather than the BERTScore variant for short-answer tasks.

\subsection{VC Degeneracy as Meta-Cognitive Failure}

The ECE = 0.430 and 5-value degenerate distribution confirm that 7B-scale instruction-tuned models lack reliable meta-cognitive access to uncertainty.
This independently replicates \citet{Xiong2023} with a different elicitation prompt, confirming the finding is robust to prompt variation.

\subsection{Limitations}

\textbf{L1 (Task-structure scope):} SE $>$ TE holds on TriviaQA but not TruthfulQA.
We frame the task-structure condition as a contribution, not a failure --- identifying \textit{when} a method works is more valuable than false universal claims.

\textbf{L2 (N=98 pilot power):} CI half-widths of ${\sim}0.05$--$0.07$ provide ample power for the +0.155 primary finding.
TruthfulQA results (AUROC range 0.44--0.51) are directional --- overlapping CIs do not definitively establish the ranking.

\textbf{L3 (NLI model size confound in H-C1):} H-C1 changes both benchmark and NLI model size simultaneously.
We cannot attribute the TruthfulQA reversal purely to task structure without the follow-up experiment.

\textbf{L4 (Sign convention inconsistency):} H-M3 and H-M4 report SE AUROC = 0.286 (uninverted).
All within-pipeline comparisons are internally consistent; corrected values are provided in Results.

\textbf{L5 (H-M2 design flaw):} The planned NLI-clustering ablation was invalidated by a monotone-transform design error.
Mechanistic evidence comes from H-M1's direct variance measurement instead.

\subsection{Broader Impact}

This work characterizes when SE is and is not a reliable uncertainty estimator, providing practitioners with empirical guidance for method selection.
The task-structure condition --- paraphrase-rich vs.\ deterministic-wrong incorrect outputs --- is a benchmark property practitioners can evaluate before committing to a method.
The VC degeneracy finding has direct safety implications: 7B-scale models expressing 95\% confidence while answering incorrectly ${\sim}53\%$ of the time should not be used as uncertainty proxies without calibration correction.
No foreseeable misuse: uncertainty estimation is defensive, and characterizing its failure modes improves safe application.
